\section{Related Work}\label{section::related_work}

This section reviews related work in four parts: the business problem in the emergency department (Section~\ref{subsection::business_problem}), the prediction of admissions and process outcomes (Section~\ref{subsection::prediction_work}), fairness in predictive process monitoring (Section~\ref{subsection::fairness_ppm}), and the resulting research gap.
% TODO: add \ref{} for the research gap once 2.4 exists.

\subsection{The Business Problem}\label{subsection::business_problem}

Patients in the emergency department (ED) are not always treated alike. At the Beth Israel Deaconess Medical Center in Boston, the hospital behind the MIMIC-IV-ED data used in this thesis, a study that has not yet been peer reviewed found that Black, Asian and Hispanic patients were less likely to be admitted to the hospital than White patients, and men were more likely to be admitted than women~\cite{coggan2025demographic}. These differences remained after the authors adjusted for the triage urgency assigned, the vital signs during the visit and the severity of the final diagnosis. Differences can start as early as triage, the first assessment of how urgent a patient's case is. In an urban academic ED in the northeastern US, White patients received more urgent triage scores than Black and Hispanic patients with the same symptom~\cite{joseph2022association}.

Insurance is another line along which treatment differs. In a national US sample of 215,028 ED visits for pneumonia and two chronic lung diseases, uninsured patients were more often sent home or transferred to another hospital instead of being admitted than privately insured patients, and Medicaid patients were more often transferred~\cite{venkatesh2019association}. This held after adjusting for how sick patients were, and all included hospitals were able to treat these patients. The authors see this as a sign that financial incentives may affect access to hospital care. In another US hospital, admitted Medicaid patients waited about 85 minutes longer for an inpatient bed than patients with commercial insurance, even though beds were assigned by a standard procedure~\cite{prucnal2025patterns}. In the MIMIC-IV-ED event log, insurance and language were the attributes most strongly associated with the discharge decision~\cite{andreswari2026fairness}.

The evidence on insurance comes mainly from later decisions in the visit, such as the admission decision and the wait for a bed. For triage, the sources reviewed here say little; one study mentions in a single sentence that no differences were seen by insurance status or language~\cite{joseph2022association}. For prediction, there is a practical problem: insurance is often unknown early in a visit. In MIMIC-IV, insurance is recorded in the table of hospital admissions, while the table of ED visits has no insurance field~\cite{mimicdocs}. Insurance is therefore missing for visits that do not end in a hospital admission. This may also affect the link reported by Andreswari et al., because visits without an admission lack insurance data. For patients who were in the hospital before, however, insurance is already on record from the earlier stay. Admission prediction models already use earlier care contacts as predictors~\cite{brink2022predicting}, and in one deployed model a previous admission was among the most important features early in the visit~\cite{king2022machine}. Historical data therefore makes it possible to examine insurance from the first decision point of a visit. This has a price. It covers only returning patients, and their insurance may have changed since the last stay; Medicaid coverage, for example, is often lost for short periods~\cite{venkatesh2019association}. The method section discusses these limits.
% TODO: replace "The method section" by Section~\ref{...} once the method section exists.

Hospitals have a business reason to predict admissions early. ED crowding leads to poorer patient outcomes and makes it harder for staff to follow recommended treatment~\cite{morley2018emergency}. A major cause of crowding is access block: admitted patients cannot move from the ED to an inpatient bed once their treatment in the ED is finished~\cite{morley2018emergency}. They stay in the ED and wait, which is called boarding. Longer boarding is associated with higher hospital mortality and longer hospital stays~\cite{singer2011association}. If a hospital knows early which patients will need a bed, it can request beds sooner. A simulation study based on data from one large ED, not yet peer reviewed, suggests that this could shorten boarding by 30 to 70 percent~\cite{cheng2026proactive}.

A review of 16 admission prediction models found that all of them used information available at triage, and that none had been put into practice~\cite{brink2022predicting}. One exception is a tool at University College London Hospitals. It predicts each patient's probability of admission at twelve moments of the visit, from arrival up to twelve hours later, and sends the expected number of admissions to bed planners four times a day~\cite{king2022machine}. Its accuracy depended on how long a patient had already been in the ED. Before the Covid-19 pandemic, the AUROC, a common measure of how well a model separates admitted from non-admitted patients, ranged from 0.82 to 0.90 across these moments. The moment of prediction is therefore a choice with consequences. Earlier predictions are usually less accurate, but they leave more time to act~\cite{metzger2023automatically}. In predictive process monitoring, alarm systems whose thresholds change as the case progresses lowered the net cost of undesired outcomes~\cite{fahrenkrogpetersen2022fire}. Staff also act on predictions that are updated as new events arrive, so a prediction that keeps changing can mislead them~\cite{teinemaa2018temporal}.

Why should differences by insurance count as unfair? Insurance is not a medical criterion. It records how care is paid for, not what care a patient needs. Pot et al.\ distinguish inequality, a mere difference between groups, from inequity, a difference that is unfair or stems from injustice; their examples of inequity include unequal access to healthcare and health differences that follow income~\cite{pot2021not}. US emergency care was shaped by this concern. The Emergency Medical Treatment and Active Labor Act (EMTALA) was passed after hospitals had refused care to uninsured patients or transferred them for financial reasons~\cite{venkatesh2019association}. It requires hospitals to examine anyone who comes to the ED and forbids delaying this to ask about payment or insurance~\cite{emtala}. Insurance also reflects income: Medicaid covers people with very low income, and R\"o\"osli et al.\ use payer type as a measure of socioeconomic status~\cite{roosli2022peeking}. It is linked to other attributes as well. Patients with limited English were far more often covered by Medicaid, safety-net insurance or self-pay than other patients (41 vs.\ 15 percent)~\cite{schulson2018emergency}. A model can therefore use insurance as a stand-in for protected attributes. Such proxy discrimination is not covered by the audit duty of the current US rule~\cite{lacava2026future}. Not every difference by insurance is unfair, however. Insurance is linked to age, because Medicare covers people over 65~\cite{roosli2022peeking}, and patients with different insurance may differ in medical need in ways the data do not show. Fairness therefore has to be judged among patients with the same medical need.

Some group differences can even be justified. A patient's sex and medical history, for example, are often essential when a physician chooses a treatment~\cite{lepsien2026role}. Past decisions are no neutral benchmark either: event logs may contain systemic inequalities, which become legitimised when they are treated as ground truth~\cite{revoredo2026towards}. Which fairness measure fits a decision is a choice that practitioners have to make~\cite{feuerriegel2020fair}. De-Arteaga et al.\ argue that the common assumption of a trade-off between fairness and utility is often mistaken, for example when the outcome a model learns from is mismeasured~\cite{dearteaga2022algorithmic}.

US law adds pressure. Since May 2025, health care providers that receive federal funding must make reasonable efforts to identify patient care decision support tools that use race, color, national origin, sex, age or disability as input, and to reduce the risk of discrimination from these tools (45~CFR \S~92.210)~\cite{hhs2024nondiscrimination,lacava2026future}. Insurance is not on this list. EMTALA requires stabilising treatment or an appropriate transfer for emergency conditions, but it does not say who must be admitted~\cite{emtala}. Differences by insurance in later decisions are thus left largely to the hospital's own judgement.

In short, decisions in the ED already differ between patient groups, including by insurance. Hospitals want admission predictions at specific moments of the visit so that they can act on them. A prediction tool trained on past decisions can carry existing differences into exactly these moments. Table~\ref{table::business_problem} summarises the evidence. The next subsection reviews how prediction research has dealt with this.

\begin{table}[htbp]
  \centering
  \footnotesize
  \caption{Evidence for the business problem}\label{table::business_problem}
  \begin{tabularx}{\textwidth}{@{}>{\raggedright\arraybackslash}p{2.6cm}>{\raggedright\arraybackslash}p{2.0cm}>{\raggedright\arraybackslash}X>{\raggedright\arraybackslash}X@{}}
    \toprule
    \textbf{Topic} & \textbf{Sources} & \textbf{What the sources show} & \textbf{What it means for this thesis} \\
    \midrule
    \multicolumn{4}{@{}l}{\textit{Is there unequal treatment in the ED?}} \\
    Triage and admission &
      \cite{joseph2022association,coggan2025demographic} &
      Triage scores and admission differ by race and sex, also after adjusting for how sick patients are. &
      The admission label in MIMIC-IV-ED may already carry group differences. \\
    Insurance &
      \cite{venkatesh2019association,prucnal2025patterns} &
      Uninsured and Medicaid patients are more often transferred or sent home; Medicaid patients wait longer for a bed. &
      Insurance differences appear mainly in later decisions of the visit. \\
    Insurance in MIMIC &
      \cite{andreswari2026fairness,mimicdocs} &
      Insurance is strongly linked to the discharge decision, but is recorded only for admitted visits. &
      Insurance must come from earlier stays; the study covers returning patients only. \\
    \midrule
    \multicolumn{4}{@{}l}{\textit{Why predict early, and why does the moment matter?}} \\
    Crowding and boarding &
      \cite{morley2018emergency,singer2011association,cheng2026proactive} &
      Crowding harms patients; long boarding goes with higher mortality; early bed requests could shorten boarding (simulation). &
      Early admission predictions have business value. \\
    Admission prediction in practice &
      \cite{brink2022predicting,king2022machine} &
      Most models use triage information only and are not in use; one deployed tool predicts at twelve moments for bed planning. &
      Predictions at several moments are realistic; their fairness has not been studied. \\
    Moment of acting &
      \cite{metzger2023automatically,fahrenkrogpetersen2022fire,teinemaa2018temporal} &
      Earlier predictions are less accurate but leave more time to act; thresholds and stability matter for decisions. &
      The decision point is a business choice and should be part of the evaluation. \\
    \midrule
    \multicolumn{4}{@{}l}{\textit{When is a difference unfair, and what does the law say?}} \\
    Inequality and inequity &
      \cite{pot2021not,lepsien2026role,revoredo2026towards} &
      A difference is inequitable if it is unfair; some attributes are legitimate in treatment; past decisions are not neutral. &
      Fairness must be judged among patients with the same medical need. \\
    Insurance as a proxy &
      \cite{roosli2022peeking,schulson2018emergency,lacava2026future} &
      Insurance reflects income and is linked to language; proxies are not covered by the audit duty. &
      Insurance is a meaningful sensitive attribute even though it is not legally protected. \\
    US law &
      \cite{hhs2024nondiscrimination,emtala} &
      Providers must check tools that use protected attributes; EMTALA protects screening, not admission. &
      Fairness by insurance in admission decisions is left to the hospital. \\
    \bottomrule
  \end{tabularx}
\end{table}

\subsection{Prediction of Admissions and Process Outcomes}\label{subsection::prediction_work}

Three lines of research come close to this problem. The first predicts hospital admission in the ED. The second, predictive process monitoring, makes predictions about processes that are still running, for example about how they will end. The third checks whether clinical prediction models treat patient groups fairly. Each line covers only part of the problem. Table~\ref{table::prediction_work} compares the studies.

Several studies predict admission on MIMIC-IV-ED. Xie et al.\ built a benchmark of 441,437 adult visits, with open code and shared outcome definitions~\cite{xie2022benchmarking}. Their hospitalisation model predicts at triage, even before a physician is assigned to the patient. Three of the studies reviewed here also check such models for fairness on the same data. Lett et al.\ compared ways to reduce errors in groups defined by race and ethnicity and by gender. Models that were adjusted to spread errors more evenly across groups, so-called debiased models, made fewer errors in several groups and kept the overall AUROC at 0.85~\cite{lett2025intersectional}. Liu et al.\ searched many almost equally accurate models for the fairest one~\cite{liu2024faim}. On MIMIC-IV-ED, the model they chose left out race and sex. Its equalized odds value, which measures the largest difference in error rates between the groups, fell from 0.316 to 0.147, while the AUC changed only from 0.790 to 0.786.

Licoppe et al.\ tested what happens when race is added as a model input~\cite{licoppe2026assessing}. The sensitivity, the share of admitted patients the model detects, dropped from 0.535 to 0.432 for Black patients and from 0.489 to 0.389 for Hispanic patients. The authors conclude that race may be better used to audit models than as a model input. All three fairness studies make a single prediction early in the visit, and none of them studies insurance or language. The tool by King et al.\ predicts at twelve moments of the visit, but it does not report fairness~\cite{king2022machine}.

Wei et al.\ turned MIMIC-IV-ED into an event log called MIMICEL~\cite{wei2023mimicel}. Each of its 425,028 ED stays is a case, that is, a sequence of events such as triage, vital sign checks, medicine dispensation and discharge. With this log, an ED visit can be studied as a process. Madau and Semeraro used MIMICEL to predict the disposition, that is, where a patient goes after the ED~\cite{madau2025process}. Their features cover the whole stay up to discharge, so the prediction comes only at the end of the visit, and they do not examine fairness. Predictive process monitoring makes predictions about such a process while it is still running. Outcome-oriented methods predict how a case will end. They train a classifier on prefixes, the first events of a case, and apply it to cases that have not finished yet. Teinemaa et al.\ compared these methods on 24 prediction tasks from nine real-life event logs and reported the AUC for each prefix length~\cite{teinemaa2019outcome}. Such results need care. As prefixes get longer, short cases have already finished, so each prefix length is evaluated on a different set of cases~\cite{teinemaa2018temporal}. A research agenda for the field describes the choice of the best moments to predict as only partly investigated, and fair decisions as a poorly explored issue for process mining as a whole~\cite{ceravolo2024predictive}. None of the benchmark studies above examines fairness.

Fairness research on clinical models shows that prediction errors can differ by insurance. In MIMIC-III, patients with Medicaid received clearly worse mortality predictions than privately insured patients, although both groups had almost the same mortality rate (8.5 vs.\ 8.3 percent)~\cite{roosli2022peeking}. In one US ED, a model that predicted long waits for an examination room missed such waits more often for men, non-Hispanic White patients and insured patients than for women, Hispanic patients and uninsured patients~\cite{wang2025evaluating}. How to reduce such differences is less clear. Pfohl et al.\ penalised differences in the predictions between groups and found that this lowered performance within almost all groups~\cite{pfohl2021empirical}. Sánchez-Marqués et al.\ found that reweighting, which gives training cases different weights, improved fairness only when the bias was low. Equalized odds post-processing, which adjusts the predictions after training, also worked under strong bias and kept performance high~\cite{sanchezmarques2026addressing}.

Which fairness measure fits depends on the decision. Rajkomar et al.\ link the measure to the harm. If missing a patient is the harm, equal sensitivity across groups matters; if false alarms lead to harmful interventions, equal specificity matters~\cite{rajkomar2018ensuring}. Paulus and Kent call a decision polar when, from the patient's point of view, one side of the prediction is always better than the other, for example when predictions allocate scarce resources. In their view, fairness concerns apply specifically to such decisions~\cite{paulus2020predictably}. An early bed request can be read as a polar decision. Paulus and Kent also warn that a feature such as previous hospitalisation can be biased if admission thresholds differed between groups in the training data. Lett et al.\ use the number of previous admissions as an input~\cite{lett2025intersectional}. A similar risk applies to this thesis. It takes insurance from earlier hospital stays, so which patients are included depends on past admission decisions.

In the studies reviewed here, fairness of ED admission prediction has been checked for race and sex, but only for one early prediction per visit. Insurance appears as a sensitive attribute in other clinical models~\cite{roosli2022peeking,wang2025evaluating}. Insurance and language have also been studied as differences in how ED patients are treated~\cite{venkatesh2019association,schulson2018emergency}, but not in admission prediction. None of these studies checks fairness at more than one moment of a visit or case. Section~\ref{subsection::fairness_ppm} turns to fairness research in predictive process monitoring, where the second and third lines of research meet.

\begin{table}[htbp]
  \centering
  \footnotesize
  \caption{Related prediction studies}\label{table::prediction_work}
  \begin{tabularx}{\textwidth}{@{}>{\raggedright\arraybackslash}p{2.9cm}>{\raggedright\arraybackslash}X>{\raggedright\arraybackslash}X>{\raggedright\arraybackslash}p{2.3cm}>{\raggedright\arraybackslash}p{2.6cm}@{}}
    \toprule
    \textbf{Study} & \textbf{Prediction task} & \textbf{Data} & \textbf{Prediction moment(s)} & \textbf{Fairness checked} \\
    \midrule
    \multicolumn{5}{@{}l}{\textit{Admission prediction in the ED}} \\
    Xie et al.~\cite{xie2022benchmarking} & Hospitalisation & MIMIC-IV-ED (441,437 visits) & Triage & No \\
    Lett et al.~\cite{lett2025intersectional} & Admission & MIMIC-IV-ED; one children's ED & Triage (children's ED: first 60 minutes) & Race and ethnicity, gender \\
    Liu et al.~\cite{liu2024faim} & Admission & MIMIC-IV-ED; one ED in Singapore & One, early in the visit (triage score, vital signs, past records) & Race, sex \\
    Licoppe et al.~\cite{licoppe2026assessing} & Admission & MIMIC-IV-ED & Triage & Race \\
    King et al.~\cite{king2022machine} & Admission, summed up for bed planning & One ED in London & Twelve moments, from arrival to 12 hours & No \\
    \midrule
    \multicolumn{5}{@{}l}{\textit{Process mining and predictive process monitoring}} \\
    Madau \& Semeraro~\cite{madau2025process} & Disposition (eight and five classes) & MIMICEL & End of the visit (data from the whole stay) & No \\
    Teinemaa et al.~\cite{teinemaa2019outcome} & Case outcome & Nine real-life event logs (24 tasks) & Each prefix length, up to 20 or 40 events & No \\
    Teinemaa et al.~\cite{teinemaa2018temporal} & Case outcome, stability of predictions & Six real-life event logs (12 tasks) & Each prefix length, up to a set limit & No \\
    \midrule
    \multicolumn{5}{@{}l}{\textit{Fairness of clinical prediction models}} \\
    Röösli et al.~\cite{roosli2022peeking} & In-hospital mortality & MIMIC-III (ICU); Stanford hospital data & First 48 hours of the ICU stay & Gender, ethnicity, insurance \\
    Wang et al.~\cite{wang2025evaluating} & Wait of 30 minutes or more from triage to an examination room & One US ED (triage level 3) & After triage & Sex, race and ethnicity, insurance \\
    Pfohl et al.~\cite{pfohl2021empirical} & Mortality, long stay, readmission & Stanford hospital data, claims, MIMIC-III & Start of the hospital stay (MIMIC-III: 24 hours after admission) & Race and ethnicity, sex, age \\
    Sánchez-Marqués et al.~\cite{sanchezmarques2026addressing} & Glioma grade & TCGA (839 patients) & Not time-based (one diagnosis per patient) & Race, gender \\
    \bottomrule
  \end{tabularx}
\end{table}

\subsection{Fairness in Predictive Process Monitoring}\label{subsection::fairness_ppm}

Fairness is a young topic in process mining. Revoredo et al.\ propose FairPM, a taxonomy of 16 biases in process mining and 12 interventions to reduce them~\cite{revoredo2026fairpm}. They describe bias as moving in a loop. Biased data shape the learned model, the model influences the decisions of its users, and these decisions become the data for the next round. Historical bias, for example, enters the loop when the data record past decisions that were already unequal. The interventions are grouped by the stage at which they act. Pre-processing removes bias from the data before training, in-processing changes the learning itself, and post-processing corrects the trained model or its output. A systematic review by Andreswari et al.\ found 42 studies on fairness in process mining and describes the field as still emerging and mainly conceptual, with little empirical work~\cite{andreswari2025review}. Twelve of the reviewed studies used real event logs, but none of these logs contained sensitive attributes such as gender, age or race. According to the review, most studies therefore added such attributes artificially.

An early fairness-aware technique in process mining addressed root cause analysis, which looks for the reasons behind problems such as delays. Qafari and van der Aalst learned a decision tree that explains such a problem and then relabelled some of its leaves, so that the tree depends less on a sensitive attribute~\cite{qafari2019fairness}. Their example of an unfair diagnosis is naming the busiest employee as the cause of a bottleneck. They also point out that removing the sensitive attribute is not enough, because rules based on other attributes can act as proxies for it. For their experiments, they added discrimination artificially to a real event log from hospital billing.

Later studies make predictions for running cases. De Leoni and Padella train a neural network together with an adversary, a second network that tries to predict the protected attributes from the output of the first network's last layer~\cite{deleoni2025achieving}. The first network is trained so that the adversary predicts them poorly. In four case studies, this reduced the contribution of the protected attributes to the predictions, measured with Shapley values, which show how much each input adds to a prediction. Accuracy stayed high. Peeperkorn and De Vos add a second term to the loss function of an LSTM model: the Wasserstein distance between the distributions of predicted probabilities in two groups~\cite{peeperkorn2025achieving}. A weight on this term sets the trade-off between predictive performance and demographic parity, the requirement that groups receive positive predictions equally often. They also measure demographic parity without a fixed decision threshold, with two metrics (ABPC and ABCC) that compare the whole distributions of predicted probabilities between the groups. Da Silva et al.\ compared one method per stage on simulated logs of a hiring process: reweighing of training cases, adversarial debiasing and equalized odds post-processing~\cite{dasilva2025towards}. Reweighing improved fairness with little loss in accuracy. Adversarial debiasing improved fairness more strongly but lowered accuracy and recall, and post-processing changed fairness little. In a clinical study, by contrast, post-processing worked better than reweighting when the bias was strong~\cite{sanchezmarques2026addressing}. Which stage works best thus seems to depend on the data.

Three of these studies use the same simulated event logs. Pohl et al.\ published twelve logs with 10,000 cases each, covering hiring, hospital treatment, lending and renting, each with a known level of discrimination~\cite{pohl2023collection}. The hospital logs also include private insurance among the patient attributes. The authors note, however, that simulated logs may not capture all the complexity of real processes. Peeperkorn and De Vos left the hospital process out because its outcome did not correlate directly with the sensitive attributes~\cite{peeperkorn2025achieving}.

Käppel et al.\ start from the observation that the same attribute can be used fairly in one decision and unfairly in another. In a medical process, for example, a treatment decision may depend on gender, while the decision to accept a patient should not~\cite{kappel2026human}. Their approach brings in human experts. It distils a decision tree from a neural network that predicts the next activity, and experts delete the splits on sensitive attributes that they judge to be unfair. The changed tree is used to relabel the training data, and the network is fine-tuned on them. The approach was tested on a simulated cancer screening process and on two real event logs, one from a Dutch financial institution and one from hospital billing. The authors added gender to both and, in the hospital billing log, also age. FairLoop offers this approach as a tool with a graphical interface~\cite{mohrlein2025fairloop}. Its demonstration uses the simulated log, \enquote{since no real life event log containing sensitive data has been made publicly available}.

De Santis et al.\ move from prediction to recommendation~\cite{desantis2027fairness}. Their method predicts a performance indicator for a running case and recommends the next activity. It uses counterfactual fairness: a recommendation should not change if only the sensitive attribute of the case is changed. The authors motivate this with a hypothetical emergency department. If older patients are sent to cardiology more often than younger ones, a quota per group could send some older patients elsewhere regardless of their symptoms. The group statistics would look fairer, but individual patients could be harmed. For testing, they cut each running case at one random point, so their results average over different moments of the cases. The sensitive attributes are again added ones, taken from Käppel et al., plus the country of an organisation in an incident management log.

Table~\ref{table::fairness_ppm} compares the six studies on predictive process monitoring. They show that changing the data or the training can reduce bias in predictions on process data. Their data, however, are far from a real care process. None of the studies uses real sensitive attributes of the people whose cases are recorded. The attributes are simulated, added by the authors, or describe staff and organisations. The logs that come from hospitals are either simulated or record billing, not patient care. Four studies train on prefixes or short windows of running cases, but none reports fairness separately for different moments of a case. Three split their data by time, so that models are trained on earlier cases and tested on later ones. Three papers also note that the recorded outcomes may themselves be biased. FairLoop measures performance on the original data, so fairer predictions can look worse, because they \enquote{conflict with the potentially biased ground truth}~\cite{mohrlein2025fairloop}. Peeperkorn and De Vos assume that the process itself is made fairer at the same time, so that a fair model should reflect these fairer executions rather than past bias~\cite{peeperkorn2025achieving}.

Among the studies reviewed here, only one examines fairness in a process with real attributes of patients. Andreswari et al.\ analysed the triage process in MIMICEL for differences by age, gender, race, language and insurance~\cite{andreswari2026fairness}. They did not build a prediction model. Instead, they asked whether the process itself treats groups equitably, and used statistical tests to compare the duration of the stay, repeated activities, deviations from a discovered process model and the discharge decision between groups. As future work, they suggest fairness-aware process monitoring tools for emergency care that could support real-time auditing.

Outside process mining, one study measured fairness along a running case, in education. Rzepka et al.\ predicted whether learners would leave a session on an online spelling platform and repeated the prediction after every sentence~\cite{rzepka2022fairness}. As the session went on, accuracy rose and fairness fell. Later predictions, however, only cover sessions that were still running, and the paper does not discuss this.

So far, fairness methods for process predictions have been tested on simulated logs, on attributes that the authors added themselves, or on attributes of staff and organisations. Among the studies reviewed here, the only one with real patient attributes from an emergency department describes the process but makes no predictions. The only one that followed fairness along running cases comes from education. The next subsection sets out the resulting gap and how this thesis addresses it.
% TODO: replace "The next subsection" by Section~\ref{...} once 2.4 exists.

\afterpage{%
\begin{landscape}
\begin{table}[p]
  \centering
  \scriptsize
  \setlength{\tabcolsep}{3pt}
  \caption{Fairness studies in predictive process monitoring compared with this thesis}\label{table::fairness_ppm}
  \makebox[\linewidth][c]{%
  \begin{tabularx}{25cm}{@{}>{\raggedright\arraybackslash}p{2.0cm}>{\raggedright\arraybackslash}p{1.8cm}>{\raggedright\arraybackslash}p{1.7cm}>{\raggedright\arraybackslash}p{2.3cm}>{\raggedright\arraybackslash}p{2.4cm}>{\raggedright\arraybackslash}p{1.8cm}>{\raggedright\arraybackslash}p{1.7cm}>{\raggedright\arraybackslash}p{1.2cm}>{\raggedright\arraybackslash}p{1.2cm}>{\raggedright\arraybackslash}p{1.3cm}>{\raggedright\arraybackslash}X>{\raggedright\arraybackslash}X@{}}
    \toprule
    \textbf{Study} & \textbf{Task} & \textbf{Synthetic log} & \textbf{Real log} & \textbf{Sensitive attribute} & \textbf{Mitigation stage} & \textbf{Running cases} & \textbf{Fairness per moment} & \textbf{Time-aware split} & \textbf{Decision named} & \textbf{Fairness metric(s)} & \textbf{Performance metric(s)} \\
    \midrule
    de Leoni \& Padella~\cite{deleoni2025achieving} & Outcome; total time & Yes (hiring, hospital) & Yes: incident management (VINST) & Simulated; real, of staff and organisations & In (adversarial) & All prefixes, pooled & No & Yes & Partly & Shapley values; equalized odds & Absolute percentage accuracy; F-score \\
    Peeperkorn \& De Vos~\cite{peeperkorn2025achieving} & Outcome & Yes (hiring, lending, renting; hospital log left out) & No & Simulated & In (Wasserstein term in the loss) & Prefixes up to 6 events, pooled & No & No & Yes & Demographic parity at a threshold; ABPC and ABCC (no threshold) & AUC, accuracy, F1 \\
    da Silva et al.~\cite{dasilva2025towards} & Outcome (job offer) & Yes (hiring) & No & Simulated & Pre, in and post & No (complete cases) & No & Yes & Partly & Disparate impact, equalized odds & Accuracy, precision, recall, F1 \\
    Käppel et al.~\cite{kappel2026human} & Next activity & Yes (cancer screening) & Yes: loan applications; hospital billing & Simulated; added (gender, age) & In (manually by an expert) & Windows of 3 events & No & No & Yes & Demographic parity & Accuracy \\
    Möhrlein et al.~\cite{mohrlein2025fairloop} & Next activity (tool) & Yes (cancer screening) & No & Simulated & In (manually by an expert) & Not described & No & Not described & Yes & None (experts judge the rules) & Accuracy, precision, recall, F1 \\
    De Santis et al.~\cite{desantis2027fairness} & Next activity to recommend & No & Yes: loan applications; hospital billing; incident management & Added (gender, age); real, of organisations & In (fairness rule during training) & One random prefix per test case & No & Yes & Yes & Counterfactual (CF-MSE, CF-RECC) & MAE; simulated effect of recommendations \\
    \midrule
    This thesis & Outcome (hospital admission) & No & Yes: MIMICEL (emergency department) & Real, of patients (insurance) & To choose & Two decision points & Yes & Yes & Yes & To choose & To choose \\
    \bottomrule
  \end{tabularx}}

  \smallskip
  \makebox[\linewidth][c]{\parbox{25cm}{\scriptsize\raggedright \textit{Sensitive attribute:} simulated = part of a simulated log; added = added to a real log by the authors; real = recorded in the real log. \textit{Running cases:} how a study handles cases that are still running. \textit{Fairness per moment:} fairness reported separately for different moments of a case. \textit{Time-aware split:} models trained on earlier cases and tested on later ones. \textit{Decision named:} the paper names the real decision that the result supports (partly: only in general terms).}}
\end{table}
\end{landscape}
}
